\section{An aggregate bound for a fixed number of ports}
\label{sec:terminal-aggregate}

The unrestricted polynomial forest estimate constrains the entire terminal-state vector of a pole.
It applies before a cut-state classification.

\begin{theorem}[All even terminal forest sectors]
\label{thm:fixed-port-aggregate}
Fix an integer $k\ge2$.
Let $T$ be a connected normalized $k$-pole of rank $h$.
For each even subset $S$ of its ports, let $F_S(T)$ count the forests whose odd-degree vertex set is exactly $S$.
Then, with $\kappa=2^{-145}$ and for all sufficiently large $h$,
\[
 \sum_{S\text{ even}}F_S(T)
 \le 2^{h+k-2}+2^{h+k-1}(h+k-1)^{-\kappa}.
\]
In particular, for a five-pole,
\[
 2^{-h}\left(\sum_{|S|=2}F_S(T)+\sum_{|S|=4}F_S(T)\right)
 \le8+16(h+4)^{-\kappa}-2^{-h}.
\]
The two-terminal sector counts are the simple terminal-path counts.
\end{theorem}

\begin{proof}
Take a tree $Q$ with $k$ leaves and degree three at every nonleaf vertex.
For $k=2$, use a single edge.
Add one edge from each leaf of $Q$ to a different port of $T$.
Call the resulting graph $\Gamma$, and put $P=V(Q)$.
Then $\Gamma$ is connected and subcubic, every vertex outside $P$ has degree three, and $\Gamma[P]$ is a tree.
Moreover,
\[
 \beta(\Gamma)=h+k-1,
 \qquad |P|=2k-2.
\]
The marked-set hypotheses of the polynomial forest theorem hold for sufficiently large $h$.

Every even subgraph of $\Gamma$ restricts to an edge set of $T$ whose odd vertices form an even port set $S$.
Conversely, any such edge set has a unique even extension to $\Gamma$.
Its port parities determine the $k$ attachment edges.
The tree edges are then uniquely determined by parity across each tree cut.
An extension exists because $|S|$ is even.
Thus each even port pattern has exactly $2^h$ preimages, and all $2^{k-1}$ patterns occur.

The restriction to $T=\Gamma-P$ is a forest exactly in the event estimated by the marked forest theorem.
Therefore
\[
 \frac{\sum_{S\text{ even}}F_S(T)}{2^{h+k-1}}
 \le\frac12+(h+k-1)^{-\kappa}.
\]
This proves the first assertion.
For $k=5$, the only sector sizes are zero, two, and four.
The zero sector contains only the empty forest, so its count is one.
The other sectors consist of one path or two vertex-disjoint paths, respectively.
\end{proof}

\begin{corollary}[Five-pole energy and nearly Boolean profiles]
\label{cor:five-pole-nearly-boolean}
Use the ten normalized path coordinates $a_e$ and the five normalized four-terminal forest totals $U_S$.
Let $z$ range over these fifteen coordinates.
Put
\[
 \xi_h=16(h+4)^{-\kappa},\qquad E(T)=\sum_z z^2.
\]
Then
\[
 E(T)\le8+\xi_h-\sum_z z(1-z).
\]
In particular, if any coordinate lies in $[\delta,1-\delta]$, for fixed $0<\delta<1/2$, then
\[
 E(T)\le8-\delta(1-\delta)/2
\]
for all sufficiently large $h$.
Any sequence of poles whose ranks tend to infinity and whose energies tend to at least eight therefore has all its coordinates tending to $\{0,1\}$, with exactly eight coordinates tending to one after passage to a subsequence.
Every path coordinate tending to one is eventually exactly one, by the terminal-path saturation gap.
\end{corollary}

\begin{proof}
Every coordinate lies in $[0,1]$.
Subtract $\sum_z z(1-z)$ from the aggregate bound and discard the negative term $-2^{-h}$.
The fixed-gap claim follows as $\xi_h\to0$.
For the sequence assertion, take a convergent subsequence of the fifteen coordinates in the compact cube $[0,1]^{15}$.
The limit has sum at most eight and squared sum at least eight.
Thus every limiting coordinate is zero or one, and exactly eight are one.
The universal path gap makes any path coordinate converging to one eventually saturated.
\end{proof}

This reduction is qualitative at the profile-classification stage.
The polynomial aggregate error is much larger than $h2^{-h/2}$.
A sharp-order cut theorem still needs exact structure or pairing cancellation in every nearly Boolean exceptional profile.
